\appendix
\section{Restriction geometry and elementary group lemmas}\label{app:auxiliaries}
This appendix records useful auxiliary statements about the parameter spaces and finite representations. The full image theorems above already imply irreducibility, their special linear lower bounds, and their determinant consequences. The statements below give additional restriction and generating-set information without changing those theorems.

For comparison, the symmetric-group module at an arrangement has a classical description. Let $P=\mathbf F_p^d$, $u=\sum e_i$, $U=\ker(\sum:P\to\mathbf F_p)$, $W=P/\langle u\rangle$, and let $D$ be the image of $U$ in $W$. For odd $p$ and $1\le r\le d-2$, the module $\bigwedge^rW$ is simple when $p\nmid d$. When $p\mid d$, its unique proper nonzero submodule is $\bigwedge^rD$, and its simple quotient is $\bigwedge^{r-1}D$. This is the dual of the augmentation-hook description in \citet[Theorem~2, p.~142; Propositions~1 and~4, pp.~136--138]{Peel}. In particular, no semisimplicity of the entire symmetric-group representation is asserted in the divisible case.

\subsection{Three modular group lemmas}\label{app:modular}
These elementary arguments describe cocycles, normal transvection blocks and a reflection obstruction. They require no classification of finite linear groups.

\begin{lemma}\label{lem:sl3-cocycles}
For an odd prime $p$ and the natural module $Q=\mathbf F_p^3$,
\[
H^1(\mathrm{SL}_3(\mathbf F_p),Q)
=H^1(\mathrm{SL}_3(\mathbf F_p),Q^*)=0.
\]
\end{lemma}
\begin{proof}
Let $c$ be a cocycle, so $c(gh)=c(g)+gc(h)$. Let $D_0$ be the diagonal sign subgroup of order four. Averaging over $D_0$, whose order is invertible, permits subtraction of a coboundary so that $c$ vanishes on that subgroup. More explicitly, if $v=|D_0|^{-1}\sum_{s\in D_0}c(s)$, then $c(t)=v-tv$ for $t\in D_0$; subtract this coboundary.

Write $t_{ij}(a)=1+aE_{ij}$. This element commutes with the sign matrix that is negative on coordinates $i,j$. The cocycle relation therefore makes $c(t_{ij}(a))$ fixed by that sign matrix, hence supported on the remaining coordinate. For
\[
x=t_{12}(a),\quad y=t_{23}(b),\quad
c(x)=Ae_3,\quad c(y)=Be_1,
\]
the cocycle identity and $c(g^{-1})=-g^{-1}c(g)$ give
\[
c([x,y])=-bAe_2-abAe_1.
\]
But $[x,y]=t_{13}(ab)$ has cocycle value supported on $e_2$. Taking $a,b\ne0$ forces $A=0$. Relabelling the coordinates proves vanishing on every elementary matrix. These matrices generate $\mathrm{SL}_3(\mathbf F_p)$, so the normalized cocycle is zero. The inverse-transpose automorphism identifies the dual calculation with the natural one.
\end{proof}

\begin{lemma}\label{lem:normal-transvection-blocks}
Let a finite subgroup $G\subseteq\mathrm{GL}(V)$ act irreducibly over $\mathbf F_p$ and contain a transvection. Let $N$ be its normal subgroup generated by all transvections in $G$. Then $V$ is a direct sum of pairwise nonisomorphic simple $N$-modules, permuted transitively by $G$. Each transvection has its image line in one summand and fixes every other summand pointwise.
\end{lemma}
\begin{proof}
The $N$-socle, the sum of all simple $N$-submodules, is nonzero. Normality makes it $G$-stable, so it is all of $V$. Thus restriction to $N$ is semisimple. The canonical isotypic summands are permuted transitively by $G$, since the sum over any orbit is $G$-stable.

On an isotypic summand $S^a$, every $x\in N$ acts as the direct sum of $a$ copies of its action on $S$. Its rank $\operatorname{rank}(x-1)$ there is therefore a multiple of $a$. A transvection is nontrivial and has total rank one, so on its unique nontrivial isotypic summand one has $a=1$. Transitivity gives multiplicity one on every summand. The rank-one condition then also shows that a transvection acts nontrivially on just one summand. This proves every assertion without applying Maschke's theorem.
\end{proof}

\begin{lemma}\label{lem:symplectic-reflection-obstruction}
Suppose $p$ is odd and an irreducible normal subgroup $N\triangleleft G\subseteq\mathrm{GL}(V)$ is the natural symplectic group, with $\dim V\ge3$. Then $G$ contains no rank-one reflection, meaning an element with eigenvalue $-1$ once and eigenvalue $1$ on a hyperplane.
\end{lemma}
\begin{proof}
Fix the invariant nondegenerate alternating form $B$. Its invariant alternating forms are scalar multiples of $B$. Indeed, writing another such form as $B(Lv,w)$ makes $L$ commute with $N$. Commuting with the symplectic transvection of each nonzero centre $v$ preserves its image line $\langle v\rangle$. Hence $L$ preserves every line and is scalar. Normality of $N$ now implies that each $g\in G$ is a similitude of $B$.

Suppose a reflection $R$ has positive hyperplane $H$ and multiplier $c$. If $c\ne1$, then $B|_{H\times H}=0$. But an isotropic subspace has dimension at most $\dim V/2$, whereas $\dim H=\dim V-1$; the symplectic dimension is even and at least four. If $c=1$, the negative eigenline is perpendicular to $H$, since $p$ is odd. It is also perpendicular to itself because $B$ is alternating, so it lies in the radical of $B$. Both cases are impossible.
\end{proof}

\subsection{Projection and prescribed restrictions}\label{app:projection}
The next statements preserve the full finite image while prescribing a geometric generic restriction of an equation.

\begin{lemma}\label{lem:projection-connected}
Let $k$ be algebraically closed, let $U\subset\mathbf P^N_k$ be nonempty open, and let $T\to U$ be a connected finite \textup{\'etale} cover. Let $f:\mathbf P^N\dashrightarrow\mathbf P^M$ be a linear projection whose centre $Z$ meets $U$. The inverse image in $T$ of the geometric generic fibre of $f|_{U\setminus Z}$ is connected. Consequently a finite local system on $U$ retains its full image on that geometric generic fibre.
\end{lemma}
\begin{proof}
The scheme $T$ is smooth and integral, by connectedness and normality; see \citet[Lemma~58.11.2]{StacksNormal}. Pull the blowup of $\mathbf P^N$ along $Z$ back to $T$, obtaining a smooth integral scheme $T'$ with a morphism to $\mathbf P^M$. Choose $z\in Z\cap U$ and a point of $T$ above it. Both are $k$-rational. The exceptional fibre at this point is $\mathbf P^M$, and the resolved projection maps this fibre identically to the target. Thus $T'\to\mathbf P^M$ has a section.

Put $K=k(\mathbf P^M)$ and let $L$ be the algebraic closure of $K$ in $k(T')$. This is a finite extension. Let $C$ be the normalization of $\mathbf P^M$ in $L$. It is finite, since a finite-type scheme over a field is Nagata. The morphism $T'\to\mathbf P^M$ factors through $C$: an element of $L$ integral over a coordinate ring of the target is integral over every normal local ring of $T'$ above that affine open, and hence belongs to it. These local maps glue. This is the normalization factorization of \citet[Lemma~29.54.4]{StacksNormalization}.

The exceptional section induces a section $\mathbf P^M\to C$. Its image is closed, because $C$ is finite over the target, and has the same dimension as the integral scheme $C$. It is therefore all of $C$; reducedness then makes the section an isomorphism. Consequently $L=K$.

Away from the exceptional divisor, the map is \textup{\'etale} over an open of the affine-space fibres of the linear projection. Its generic fibre is smooth. The equality $L=K$ implies geometric irreducibility of that generic fibre by \citet[Lemma~33.8.6]{StacksGeometric}. Smoothness gives geometric reducedness, hence geometric integrality. Removing the exceptional divisor leaves a nonempty geometrically integral open. This is the required inverse image in $T$.

Finally, a connected component of the finite frame cover for a local system is connected after this restriction. Its deck group is its finite image, so that image is preserved.
\end{proof}

For $j\ge2$, let $W_{j,e}$ denote the Hasse--Witt invertibility open in the full projective space of degree-$e$ equations in $\mathbf P^j$, including singular equations. Let $V_{j,e}$ be the local system of Frobenius-fixed vectors in $H^{j-1}(\mathcal O_Y)$, and let $G_{j,e}$ be its geometric image. Sections~\ref{sec:pn-1}--\ref{sec:pn-2} give its constant rank $\binom{e-1}{j}$ and its finite \textup{\'etale} construction.

\begin{proposition}\label{prop:restriction-product}
Let $k$ be algebraically closed of characteristic $p$, let $n\ge3$ and $e\ge n+1$, and fix a hyperplane $H\subset\mathbf P^n$. On the open parameter space where both a degree-$e$ equation and its restriction to $H$ have invertible Hasse--Witt maps, the joint geometric image on $V_{n,e}$ and $V_{n-1,e}$ is the product of their full geometric images. In particular, $V_{n,e}$ retains its full image when the restriction to $H$ is prescribed to be geometric generic.
\end{proposition}
\begin{proof}
Restriction of equations is a linear projection between their projective parameter spaces. Its centre consists of equations divisible by the equation of $H$. A general-position union of $e$ hyperplanes, one of them $H$, lies in that centre and in $W_{n,e}$: its intersection complex has constant coefficients and Frobenius is invertible. Hence Lemma~\ref{lem:projection-connected} applies to the finite frame cover of $V_{n,e}$.

Away from its centre, restriction is an affine-space bundle. The additional invertibility opens have nonempty intersection: the space of nonzero restrictions is dominant, and each invertibility condition is a nonempty open. Its geometric generic fibre is therefore a nonempty open of affine space, hence geometrically integral. Pullback along this map preserves connected finite \textup{\'etale} covers of the target, so the image on $V_{n-1,e}$ is full. On the geometric generic fibre, the image on $V_{n,e}$ is full by the preceding lemma and the restricted representation is constant. The joint image consequently contains $G_{n,e}\times1$ and surjects onto $G_{n-1,e}$. These two properties imply that it is the product.
\end{proof}

\begin{lemma}\label{lem:elliptic-constant-section}
Suppose $p$ is odd and $k$ is algebraically closed. There is a family of plane cubics over a strictly henselian smooth trait over $k$ whose special fibre is smooth supersingular, whose Hasse invariant has valuation one, and whose restriction to a fixed line is the constant binary form
\[
g=x_0x_1(x_0-x_1).
\]
There is also a smooth ordinary plane cubic with this exact restriction.
\end{lemma}
\begin{proof}
Section~\ref{sec:pn-5} proves that the Legendre Hasse polynomial has a simple root away from $0$ and $1$. Take the strict henselization of the parameter line at such a root. The resulting plane-cubic family is smooth, its special fibre is supersingular, and its Hasse invariant has valuation one.

Choose a line transverse to the special cubic. Its intersection with the family is finite \textup{\'etale} of degree three, so over the strictly henselian trait it splits into three disjoint sections. A regular automorphism of that projective line takes these ordered sections to $0,1,\infty$. To see regularity explicitly, represent the first two sections by primitive vectors over the local ring; disjointness makes them a basis. The third has two unit coordinates in this basis, and diagonal rescaling sends it to $(1,1)$. The resulting invertible two-by-two matrix extends to an invertible plane-coordinate matrix by adjoining one diagonal entry.

Apply this coordinate change to the entire family. Its restricted binary cubic now has the zeros of $g$ and differs from $g$ by a unit, since those three zeros remain distinct. Rescale the defining equation by the inverse unit. The restriction is exactly $g$. Coordinate changes and unit rescaling preserve smoothness, supersingularity and the valuation of the Hasse section: they change its trivialization only by a unit. The ordinary assertion follows from the prescribed-squarefree-restriction construction in Section~\ref{sec:plane-cubic-restriction}.
\end{proof}

\subsection{An explicit characteristic-two generating set}\label{app:coordinate-generators}
This calculation identifies the subgroup generated by a specific family of coordinate swaps, giving more information than their membership in the full monodromy image.

Let $n\ge2$, $d\ge n+2$, $m=d-1$, and
\[
W=\mathbf F_2^d/\mathbf F_2\mathbf1,
\qquad V=\bigwedge^nW,
\qquad h=\binom mn.
\]
Choose the classes $e_1,\ldots,e_m$ as a basis, so $e_0=\sum_{i=1}^m e_i$. Write $b_A=\bigwedge_{i\in A}e_i$ for its exterior basis and $\ell=\sum_A b_A^*$. For an $(n-1)$-subset $I\subset\{2,\ldots,m\}$ put
\[
w_I=\sum_{a\in\{2,\ldots,m\}\setminus I}b_{I\cup\{a\}},
\qquad R_I=1+w_Ib_{I\cup\{1\}}^*.
\]
The $R_I$ are nontrivial commuting involutions: their rows are distinct basis rows containing label $1$, whereas all their centres avoid that label. A product is the identity only when every chosen factor is absent, by inspection of its distinct input columns. They thus form an independent elementary abelian subgroup. In the stabilizer parity their centres have even weight, so these operators belong to the image already determined by the main theorems.

\begin{proposition}\label{prop:coordinate-generators}
Let $G_0$ be generated by the permutation action of $S_d$ on $V$ and all the $R_I$. Then
\[
G_0=\begin{cases}
\mathrm{GL}_h(\mathbf F_2),&d-n-1\text{ odd},\\
\operatorname{Stab}_{\mathrm{GL}_h(\mathbf F_2)}(\ell),&d-n-1\text{ even}.
\end{cases}
\]
In the second case the group is $\mathbf F_2^{h-1}\rtimes\mathrm{GL}_{h-1}(\mathbf F_2)$.
\end{proposition}
\begin{proof}
Set $l=d-n-1=m-n$. Conjugating by the subgroup $S_m$ fixing label $0$ gives, for each $(n-1)$-set $I$ and each $j\notin I$, the involution $T_j^I$ which sends $b_{I\cup\{j\}}$ to
\[
S_I=\sum_{a\in\{1,\ldots,m\}\setminus I}b_{I\cup\{a\}}
\]
and fixes all other basis vectors. It also sends $S_I$ back to $b_{I\cup\{j\}}$. For distinct $j,t\notin I$, the word $T_j^IT_t^IT_j^I$ therefore exchanges just $b_{I\cup\{j\}}$ and $b_{I\cup\{t\}}$. The graph of $n$-subsets joined when they differ by one label is connected: replace successively a label outside the target subset by one in the target. Its edge transpositions generate every permutation of the coordinate basis. Thus $G_0$ contains $S_h$.

Each $R_I$ has a centre of weight $l$ disjoint from its input coordinate. Coordinate conjugation now supplies every transformation $1+w b_A^*$ with any such centre. If $l=1$, these are all elementary matrices, proving the first case. Suppose $l\ge2$. Since $h=\binom mn\ge m=l+n\ge l+2$, one can choose two $l$-element supports, both avoiding a fixed input coordinate, with intersection of size $l-1$ and prescribed distinct remaining coordinates $B,C$. Multiplying the same-row transformations gives
\[
1+(b_B+b_C)b_A^*.
\]
Hence all triple operations of Lemma~\ref{lem:two-triple-stabilizer} occur.

If $l$ is odd, add weight-two centres within an $l$-element centre to cancel pairs until one coordinate remains. Same-row transformations multiply by adding their centres. Coordinate permutations then give every elementary matrix, so $G_0=\mathrm{GL}_h(\mathbf F_2)$.

If $l$ is even, the $R_I$ and the coordinate subgroup $S_m$ preserve $\ell$. The arrangement transposition $(01)$ is exactly $\prod_I R_I$: on a basis vector not containing label $1$ both sides are the identity, and on $b_{I\cup\{1\}}$ both replace $e_1$ by $e_0=\sum e_i$. Since $S_m$ and $(01)$ generate $S_d$, every generator of $G_0$ preserves $\ell$. Conversely the triple operations generate its full stabilizer by Lemma~\ref{lem:two-triple-stabilizer}, whose elementary proof includes both the linear action on $\ker\ell$ and all translations of a complementary vector. This proves the second equality and its semidirect-product description.
\end{proof}
